\chapter{Qualitative Exploration of Existing Solutions}
\label{cha:exploration}

In order to explore the potentials and limitations of existing German-language models for Named Entity Recognition, two models were selected for a basic qualitative exploration of their results with some sample texts. 

While a quantitative evaluation would be the ideal approach, this is not possible since there are no existing (German-language) models which support all the entity types relevant for the NER task as defined here. Precisely, most of the models are limited to the four common entity types PERSON, LOCATION, ORGANIZATION, and OTHER/MISC. To the best of my knowledge, there are no existing models supporting the detection of the entity types TIME, DATE, and ACTION. 
An exception is the English-language spaCy model which supports TIME and DATE recognition. Against this background, iI will explore the potentials of machine translation to allow for the use of well-resourced English NER models for the recognition of named entities in the selected example texts.

\section{German-language Named Entity Recognition with spaCy and Flair}

The models explored in this chapter were selected from the library spaCy and the framework FLAIR. The models were chosen due to their easy applicability, allowing for fast exploration and analysis, and because both of them support German- and English-language named entity recognition. 

\subsection{The German-language spaCy model}
spaCy is a free open-source library for Natural Language Processing developed for use in production. The library is extensively documented and thus serves as an easy starting point for beginner, while also claiming to be designed specifically for use in production.
The library supports more than 66 languages (including German), comes with pre-trained word vectors, and supports a variety of NLP tasks such as part-of-speech tagging, text classification, and named entity recognition (cf. \cite{spacy_facts}). 
The German spaCy model supports to the common entity types are \verb|LOC|, \verb|MISC|, \verb|ORG|, and \verb|PER| (cf. \cite{spacy_german_md}). 

The central component of the model's NER pipeline is spaCy's Entity Recognizer which `identifies non-overlapping labelled spans of tokens' (\cite{spacy_entity_recognizer}) and uses a transition-based algorithm which assumes that the most relevant information about an entity is close to its first token (cf. \cite{spacy_entity_recognizer}). The underlying model is a neural network state prediction model, spaCy's TransitionBasedParser. Transition-based parsing can be described as `an approach to structured prediction where the task of predicting the structure is mapped to a series of state transitions.' (\cite{spacy_transition_based_parser}). The model is composed of two to three subnetworks which take care of token-vectorization, feature-specific vectorization for each token-feature pair and then computing the state representation from them, and (optionally) a subnetwork which predicts scores from the state representation (otherwise the outputs from the second subnetwork will be used for prediction) (cf. \cite{spacy_transition_based_parser}).

The model was trained on the WIKINER corpus, a free, multilingual corpus created from Wikipedia articles which were automatically annotated and thus are silver standard training data (cf. \cite{nothman_learning_2013}). According to the official documentation, the model achieved an F1 score of 0.84 (cf. \cite{spacy_german_md}).

\subsection{The German-language Standard NER Flair Model}
\label{sec:flair_de}
The FLAIR framework is an NLP library released in 2018, developed to `facilitate training and distribution of
state-of-the-art sequence labeling, text classification and language models' (\cite{akbik2019flair}). The framework's aim is to provide an interface to easily use different types of word and document embeddings (cf. \cite{akbik2019flair}).

The German-language FLAIR model for NER uses a combination of GloVe embeddings (cf. \cite{pennington2014glove} and Flair embeddings to create the word embeddings. The conceptual approach behind the Flair embeddings is described in detail in \cite{akbik2018coling}. It can briefly be characterized as an approach which creates `contextual string embeddings' (\cite{akbik2018coling}) which take into account the context of a given text, resulting in different embeddings for the same word depending on its context in a document (cf. \cite{akbik2018coling}).
The word embeddings passed to a BiLSTM-CRF sequence labeling model which takes care of the final downstream sequence labeling task (cf. \cite{akbik2018coling}). 

The German Flair model used here is the German standard 4-class NER model by Flair. The model was trained on the ConLL 2003 corpus, and reached an F1 score of 87.94 according to the model's HuggingFace documentation (cf. \cite{flair_ner_german_standard_model}). 
 
\subsection{Example Data}
For the exploration of the two models, ten example messages were randomly selected from the Telegram channel `Demotermine' (\textit{Demo dates}). 

The channel was selected because a prior qualitative exploration of sample messages indicated that its content consists mainly of information about and mobilization for demonstrations and others forms of protests against the German COVID-19 containment measures and could thus be expected to contain a lot of relevant entities. 

\subsection{Implementation}
\label{sec:impl_ner_de}
For exploring the German medium-size spaCy model, a simple code was implemented to predict and print the detected entities: 

\begin{python}
import spacy
nlp = spacy.load("de_core_news_md")

def get_entities(texts): 
    for text in texts:                 
        doc = nlp(text)
        if doc.ents:
            print("Text: ", text, '\n')
            print("The following NER tags are found: ", 
            [(ent.text, ent.label_) for ent in doc.ents], 
            '\n\n')
get_entities(texts)
\end{python}

\newpage
The next code section shows the implementation for the Flair model:

\begin{python}
from flair.models import SequenceTagger
from flair.data import Sentence

tagger = SequenceTagger.load("flair/ner-german")

def get_entities(texts):    
    for text in texts: 
        # transform example text to Sentence
        sentence = Sentence(text)
        # use tagger to predict NER tags in sentence
        tagger.predict(sentence)
        print(sentence, '\n')
        print('The following NER tags are found:')
        for entity in sentence.get_spans('ner'):
            print(entity)
        print('\n\n')
get_entities(texts)
\end{python}


\subsection{Results}

Table \ref{tab:detect_ents_de} shows the results of both models for the example texts: 

\begin{table}[htp]
\caption{Detected entities for German-language example messages}
\label{tab:detect_ents_de}
\centering
\begin{adjustbox}{max width=\textwidth}
\begin{tabular}{|p{0.04\linewidth}|p{0.55\linewidth}|p{0.36\linewidth}|p{0.3\linewidth}|}
\hline
\textbf{No.} & \textbf{Text} & \textbf{spaCy/de\_core\_news\_md} & \textbf{flair/ner-german} \\
\hline 
1 
& Freiburg Aufzug mit Endkundgebung am Platz der alten
Synagoge . Live Musik ! Details 19062021 Raus auf die Straße 
& Endkundgebung[LOC], alten Synagoge[LOC]
& Freiburg[LOC] \\
\hline
2 
& GEHT DEMONSTRIEREN ! - Message von Elijah Tee an die Dresdner \& anderen Demonstranten in Deutschland Morgen am 13.06. um 16.30 Uhr und bundesweit überall , wo es Mitdenker gibt . - Bilder und Videos hochladen - Diskussion Fahren oder Mitfahren
& GEHT DEMONSTRIEREN [ORG], Elijah Tee[PER], Dresdner \&[LOC], Deutschland[LOC], Videos[MISC]
& Elijah Tee[PER], Deutschland[LOC] \\ 
\hline
3 
& 15831 Blankenfelde Mahlow Berliner Speckgürtel Ausfahrt Lichtenrade am 14.03.2021 14032021 Raus auf die Straße
& Blankenfelde Mahlow[PER], Berliner[MISC]
& Blankenfelde Mahlow[LOC], Lichtenrade[LOC]\\ 
\hline
4 
& 2021 10 11 Sven Liebich fragt Ministerpräsident Haseloff nach dessen Verständnis 
von Versammlungsfre - Sven Liebich 0:23 Raus auf die Straßen  \"Ubersicht / Overview 
& Sven Liebich[PER], Haseloff[PER], Versammlungsfre[PER], Sven Liebich[PER]] 
& Sven Liebich[PER, Haseloff[PER], Sven Liebich[PER]\\ 
\hline
5 
& F\"UR ALLE DEMOKRATEN   GEFAHR : Die Polizei blockiert den Stern !   Deswegen laufen die Menschen vom Ernst-Reuther-Platz nach Südosten zum Breitscheitplatz und dort weiter auf den Kudamm ! ! Schließt euch an : Neue Kantstraße direkt Richtung Osten auf den Breitscheitplatz ! Von der Straße des 17. Juni Richtung Süden zum Kudamm ! 
& DEMOKRATEN[ORG], Ernst-Reuther-Platz[LOC], Breitscheitplatz[LOC], Kudamm ! ![LOC]), Neue Kantstraße[LOC], Breitscheitplatz[LOC], Kudamm ![LOC]
& FU[ORG], Ernst-Reuther-Platz[LOC], Breitscheitplatz[LOC], Kudamm[LOC], Kantstraße[LOC], Breitscheitplatz[LOC], Kudamm[LOC]\\ 
\hline
6 
& Dr. Gunter Frank : Wie Moralismus dazu führt das man Feinde sieht und vernichten will — Antihygienista - Offene Gesellschaft Kurpfalz 15:04 Raus auf die Straßen   \"Ubersicht / Overview 
& Gunter Frank[PER], Antihygienista[LOC], Offene Gesellschaft Kurpfalz[ORG] 
& Gunter Frank[PER], Antihygienista[ORG], Offene Gesellschaft Kurpfalz[ORG]\\ 
\hline
7 
& Einfach nur Krank solche Politiker ! Was sollen wir mit ihm machen wenn unsere Kinder anhand dieser Impfung nen schaden erleiden ? Wer ist dann dafür verantwortlich , wen schmeißen wir dann in den Knast und nehmen sein ganzes Geld weg , sollen wir das mit dir machen ? Raus auf die Straße 
& Einfach [MISC]
& \textit{no entities detected}\\ 
\hline
8 
& 27.11.2021 Frankfurt Durchsagen von gemischt mit Musik . Super Idee   Für mehr Berichte folgen auf ... Raus auf die Straßen   \"Ubersicht / Overview 
& Frankfurt[LOC]
& Frankfurt[LOC]\\ 
\hline
9 
& Demo-Berlin : Der Mann mit dem Schild wurde soeben von der Polizei festgenommen . Circa 7 Einsatzwägen sind vor Ort und umzäunen den Park großräumig . Am Humboldthain gibt es nichts mehr zu gewinnen . \# b2908 Raus auf die Straßen 
& Der Mann mit dem Schild[LOC], Einsatzwägen[LOC], Humboldthain[LOC]
& Demo-Berlin[LOC], Humboldthain[LOC]\\ 
\hline
10 
& Anthony Fauci : Definition von `` vollständig geimpft '' könnte geändert werden — Epoch Times Nachrichten 6:59 Raus auf die Straßen   \"Ubersicht / Overview 
& Anthony Fauci[PER], Times Nachrichten[MISC]
& Anthony Fauci[PER], Epoch Times Nachrichten[ORG]

\end{tabular}
\end{adjustbox}
\end{table}

I would like to emphasize that due to the small size of the explored sample, these explorations are rather anecdotal and cannot be generalized without further quantitative evaluation on a larger sample. 

However, the following observations can be made when looking at the results: 
 
The spaCy model produces more false positives overall: It falsely tags the sequences `Endkundgebung' (final rally), `GEHT DEMONSTRIEREN' (go demonstrate), `Dresdner \&', `Videos', `Berliner', `Versammlungsfre' (freedom of assembly, truncated), `Einfach' (simple/simply), `Der Mann mit dem Schild' (the man with the sign), and `Einsatzwägen' (patrol cars) as entities. 

The Flair model only falsely tags `FU' and `Demo-Berlin'. The sequence `FU' was originally part of the German word `F\"UR' and apparently was modified by the model's tokenizer which is applied when creating a Sentence object from the input text. The sequence `Demo-Berlin' contains a location, which makes the model's mistake to some extent plausible. 

It can also be observed that the spaCy model sometimes assigns the wrong tag to a correctly detected entity. This applies for the sequences `Blankenfelde Mahlow' which is not a person as suggested by the model, but a location. 

Regarding false negatives, it can be observed that the sapCy model does not detect `Freiburg' in the first message. On the other hand, it recognizes parts of the sequence `Platz der alten Synagoge` (square of the old synagogue) as a location. The Flair model correctly detects Freiburg, but misses the mentioned square. 

Even though the sample is too small for generalization, this first exploration casts strong doubts on the claim that spaCy offers production-ready components for NER.

It can be hypothesized that both the different model architectures and the different corpora used for training result in the better performance of the Flair model compared to the spaCy model.


\section{Exploring the Potentials of Machine Translation for NER}

Even though the spaCy model explored in the previous section provided rather dissatisfying results, its English equivalent has one advantage, since it supports more entity types, including DATE and TIME, meaning that it might be useful for extracting dates and times, as they are relevant for the aim of this thesis. 
Also, since more training data is available for English-language NLP models, exploring the potentials of machine translation for NER seems worth a try. 

\subsection{The English-language spaCy Model}

The English-language spacy model supports 18 entity types: (\verb|CARDINAL|, \verb|DATE|, \verb|EVENT|, \verb|FAC|, \verb|GPE|, \verb|LANGUAGE|, \verb|LAW|, \verb|LOC|, \verb|MONEY|, \verb|NORP|, \verb|ORDINAL|, \verb|ORG|, \verb|PERCENT|, \verb|PERSON|, \verb|PRODUCT|, \verb|QUANTITY|, \verb|TIME|, and \verb|WORK_OF_ART|). 

Table \ref{tab:spacy_ents_def} shows the definition for each entity type (obtained for each entity via the command \verb|spacy.explain{'<ENTITY>'}|: 

\begin{table}[htp]
\caption{Overview of spaCy's entity types}
\label{tab:spacy_ents_def}
\centering
\begin{adjustbox}{max width=\textwidth}
\begin{tabular}{|p{0.25\linewidth}|p{0.75\linewidth}|}
\hline
\textbf{Entity type} & \textbf{Definition} \\
\hline 
CARDINAL & Numerals that do not fall under another type. \\
\hline 
DATE & Absolute or relative dates or periods. \\
\hline 
EVENT & Named hurricanes, battles, wars, sports events, etc. \\
\hline 
FAC & Buildings, airports, highways, bridges, etc. \\
\hline 
GPE & Countries, cities, states. \\
\hline 
LANGUAGE & Any named language. \\
\hline 
LAW & Named documents made into laws. \\
\hline 
LOC & Non-GPE locations, mountain ranges, bodies of water. \\
\hline 
MONEY & Monetary values, including unit. \\
\hline  
NORP & Nationalities or religious or political groups. \\
\hline 
ORDINAL & "first", "second", etc. \\
\hline 
ORG & Companies, agencies, institutions, etc. \\
\hline 
PERCENT & Percentage, including "\%".  \\
\hline 
PERSON & People, including fictional. \\
\hline 
PRODUCT & Objects, vehicles, foods, etc. (not services). \\
\hline 
QUANTITY & Measurements, as of weight or distance. \\
\hline 
TIME & Times smaller than a day \\
\hline 
WORK\_OF\_ART & Titles of books, songs, etc.\\
\hline 
\end{tabular}
\end{adjustbox}
\end{table}

The training data sources listed in the model's official documentation suggest that the NER training data was the OntoNotes 5 corpus (cf. \cite{ontonotes5}). The model achieved an F1 score of 0.85 for NER. 

\subsection{The English-language Standard NER Flair Model}
The architecture of the standard English-language Flair model is the same as for the standard German model described in section \ref{sec:flair_de}. It was trained on a corrected version of the CoNLL 2003 corpus and achieved an F1 score of 93.06 (cf. \cite{flair_ner}). 

\subsection{Implementation}

To explore the potentials of the medium-size English spaCy model and the Flair English-language standard model, the example texts were translated into English using the DeepL API: 

\begin{python}
import deepl
from src.config import deepl_credentials
auth_key = deepl_credentials.auth_key

# initialize and authenticate translator
translator = deepl.Translator(auth_key)

# method to get translated texts
def translate(text, translator, target_lang, formality="less"): 
    result = translator.translate_text(text, target_lang=target_lang)
    return result.text

# apply translation to cleaned texts
df['text_EN'] = df.cleaned_text.apply(lambda x: translate(x, translator, 'EN-US'))
\end{python}

The translated texts were then fed to the English-language spaCy and Flair model, using the same code as described in section \ref{sec:impl_ner_de}.

\subsection{Results}

Table \ref{tab:detect_ents_en} shows the results of both models for the translated example texts. 

\begin{table}[htp]
\caption{Detected entities for English-language example messages}
\label{tab:detect_ents_en}
\centering
\begin{adjustbox}{max width=\textwidth}
\begin{tabular}{|p{0.04\linewidth}|p{0.55\linewidth}|p{0.36\linewidth}|p{0.3\linewidth}|}
\hline
\textbf{No.} & \textbf{Text} & \textbf{spaCy/en\_core\_web\_md} & \textbf{flair/ner-english} \\
\hline 
1 
& Freiburg procession with final rally at the square of the old synagogue . Live music ! Details 19062021 Out on the street
& \textit{no entities detected}
& Freiburg[LOC]\\
\hline
2 
& GO DEMONSTRATE ! - Message from Elijah Tee to the Dresdeners \& other demonstrators in Germany Tomorrow on 13.06. at 16.30 clock and nationwide everywhere , where there are fellow thinkers . - Upload pictures and videos - discussion drive or ride along 
& Germany[GPE], Tomorrow[DATE], 13.06[CARDINAL], 16.30[CARDINAL]
& Elijah Tee[PER], Dresdeners[MISC], Germany[LOC]\\ 
\hline
3 
& 15831 Blankenfelde Mahlow Berlin Speckgürtel Lichtenrade exit on 03/14/2021 14032021 Out on the road
& 15831[DATE], Blankenfelde Mahlow Berlin[PERSON], Speckgürtel Lichtenrade[PERSON], 03/14/2021 14032021[DATE] 
& "Blankenfelde Mahlow Berlin Speckgürtel Lichtenrade"[MISC]\\ 
\hline

4 
& 2021 10 11 Sven Liebich asks Prime Minister Haseloff about his understanding of assembly free - Sven Liebich 0:23 Raus auf die Straßen Overview / Overview 
& Sven Liebich[ORG], Haseloff[PER], Straßen Overview / Overview[ORG]
& Sven Liebich[PER], Haseloff[PER], Sven Liebich[PER], Raus auf die Straßen Overview[ORG]\\ 
\hline
5 
& FOR ALL DEMOCRATS DANGER : The police blocks the Stern !   That's why people run from Ernst-Reuther-Platz to the southeast to Breitscheitplatz and there further to Kudamm ! ! Join : Neue Kantstraße directly to the east to Breitscheitplatz ! From the Straße des 17. Juni southwards to the Kudamm ! 
& DEMOCRATS[NORP], Ernst-Reuther-Platz[ORG], Breitscheitplatz[PERSON], Kudamm[FAC], Neue Kantstraße[PERSON], Straße[GPE], 17[CARDINAL], Kudamm[FAC] 
& Stern ![MISC], Ernst-Reuther-Platz[LOC], Breitscheitplatz[LOC], Kudamm[LOC], Neue Kantstraße[ORG], Breitscheitplatz[LOC, ]Straße des[LOC], Juni[LOC], Kudamm[LOC]\\ 
\hline

6 
& Dr. Gunter Frank : How moralism leads to the fact that one sees enemies and wants to destroy - Antihygienista - Offene Gesellschaft Kurpfalz 15:04 Raus auf die Straßen \"Ubersicht / Overview 
& Gunter Frank[PER], 15:04 Raus auf[PER], Straßen Übersicht / Overview[ORG] 
& Gunter Frank[PER]\\ 
\hline

7 
&  "Just sick such politicians ! What should we do with him if our children suffer from this vaccination ? Who is responsible for it , who do we throw in jail and take away all his money , shall we do that with you ? Out on the street" 
& \textit{no entities detected}
& \textit{no entities detected}\\ 
\hline

8 
& 27.11.2021 Frankfurt announcements of mixed with music . Super idea For more reports follow on ... Out on the streets Overview 
& Frankfurt[GPE]
& Frankfurt[LOC]\\ 
\hline

9 
& Demo-Berlin : The man with the sign has just been arrested by the police . About 7 emergency vehicles are on site and fence the park extensively . At the Humboldthain there is nothing more to win . \# b2908 Out on the streets 
 
& Humboldthain[PER], \#[CARDINAL]
& Humboldthain[LOC]\\ 
\hline

10 
& Anthony Fauci : Definition of " fully vaccinated " could be changed - Epoch Times News 6:59 Out on the Streets Overview / Overview. 

& Anthony Fauci[PER], the Streets Overview / Overview[ORG]
& Anthony Fauci[PER], Epoch Times News[ORG], Streets Overview[MISC]
\end{tabular}
\end{adjustbox}
\end{table}



\section{Summarization}

potentials and limitations / observed errors
- spacy
- Flair
explorative evaluation

